% Part: first-order-logic
% Chapter: proof-systems
% Section: axiomatic-deduction

\documentclass[../../../include/open-logic-section]{subfiles}

\begin{document}

\iftag{FOL}
      {\olfileid[fr]{fol}{prf}{axd}}
      {\olfileid[fr]{pl}{prf}{axd}}

\olsection{\usetoken{P}{derivation} axiomatiques}

Les systèmes de !!{derivation} axiomatiques sont les plus anciens et
les plus simples des systèmes logiques de !!{derivation}. Leurs
!!{derivation}s sont de simples suites d'!!{sentence}s. Une telle
suite est une !!{derivation} correcte si chacun de ses !!{sentence}s
$!A$ satisfait l'une des conditions suivantes :
\begin{enumerate}
\item $!A$ est un axiome ; ou
\item $!A$ est un !!{element} d'un ensemble donné~$\Gamma$ d'%
  !!{sentence}s ; ou
\item $!A$ est justifié par une règle d'inférence.
\end{enumerate}
Pour être un axiome, $!A$ doit avoir la forme de l'un de plusieurs
schémas fixés d'!!{sentence}s. De nombreux ensembles de schémas
d'axiomes donnent un système de !!{derivation} satisfaisant, c'est-à-dire
correct et complet, pour la logique du premier ordre. Certains sont
organisés selon les connecteurs qu'ils régissent. Ainsi, les schémas
\[
!A \lif (!B \lif !A) \qquad !B \lif (!B \lor !C) \qquad (!B \land !C) \lif !B
\]
sont des axiomes usuels pour $\lif$, $\lor$ et~$\land$.
Certains systèmes axiomatiques cherchent à réduire au minimum le nombre
d'axiomes. Selon les connecteurs choisis comme primitifs, il est même
possible de trouver des systèmes qui ne comportent qu'un seul axiome.

Une règle d'inférence est un énoncé conditionnel qui donne une condition
suffisante pour qu'un !!{sentence} d'une !!{derivation} soit justifié.
Le modus ponens est une règle très courante : si $!A$ et $!A \lif !B$
sont déjà justifiés, alors $!B$ l'est aussi. Autrement dit, une ligne
d'une !!{derivation} contenant l'!!{sentence}~$!B$ est justifiée
lorsque $!A$ et $!A \lif !B$, pour un certain !!{sentence}~$!A$,
figurent tous deux avant~$!B$ dans cette !!{derivation}.

La relation $\Proves$ fondée sur les !!{derivation}s axiomatiques se
définit ainsi : $\Gamma \Proves !A$ si et seulement s'il existe une
!!{derivation} dont la dernière formule est l'!!{sentence}~$!A$,
tous les !!{sentence}s justifiés par la condition~(2) ci-dessus étant
pris dans~$\Gamma$. $!A$ est un théorème s'il admet une
!!{derivation} où $\Gamma$ est vide, c'est-à-dire où chaque
!!{sentence} est justifié par~(1) ou par~(3). Voici, par exemple,
une !!{derivation} qui établit que
$\Proves !A \lif (!B \lif (!B \lor !A))$ :
\begin{derivation}
  1. & $!B \lif (!B \lor !A)$ \\
  2. & $(!B \lif (!B \lor !A)) \lif (!A  \lif (!B \lif (!B \lor !A)))$\\
  3. & $!A  \lif (!B \lif (!B \lor !A))$
\end{derivation}
L'!!{sentence} de la ligne~1 a la forme de l'axiome
$!A \lif (!A \lor !B)$, les rôles de $!A$ et de $!B$ étant intervertis.
Celui de la ligne~2 a la forme de l'axiome $!A \lif (!B \lif !A)$.
Les deux lignes sont donc justifiées. La ligne~3 est justifiée par
modus ponens : si on l'abrège par $!D$, la ligne~2 a la forme
$!C \lif !D$, où $!C$ est $!B \lif (!B \lor !A)$, c'est-à-dire la
ligne~1.
\begin{editorial}
Note de l'édition française : l'original identifie ici $\Gamma$
exactement aux prémisses effectivement utilisées. Nous précisons que
ces prémisses appartiennent à $\Gamma$, conformément à la définition
formelle donnée plus loin dans l'original. Le contexte peut ainsi
contenir des hypothèses inutilisées, y compris être infini, tandis
qu'une dérivation reste finie.
\end{editorial}

Un ensemble $\Gamma$ est incohérent si $\Gamma \Proves \lfalse$.
Un système axiomatique complet démontre aussi $\lfalse \lif !A$ pour
tout~$!A$. Si $\Gamma$ est incohérent, on a donc
$\Gamma \Proves !A$ pour tout~$!A$.

Les premiers systèmes de !!{derivation} axiomatiques pour la logique
ont été donnés par Gottlob Frege dans sa \emph{Begriffsschrift} de 1879,
souvent considérée, pour cette raison, comme le premier ouvrage de
logique moderne. Ils ont été perfectionnés dans les
\emph{Principia Mathematica} d'Alfred North Whitehead et de Bertrand
Russell, puis par David Hilbert et ses élèves dans les années 1920.
On les appelle donc souvent « systèmes de Frege » ou « systèmes de
Hilbert ». Ils sont très polyvalents : il est souvent facile de trouver
un système axiomatique pour une logique donnée. Comme leurs
!!{derivation}s ont une structure très simple et ne font appel qu'à
une ou deux règles d'inférence, il est aussi relativement facile de
démontrer des résultats \emph{à leur sujet}. Ils sont cependant très
difficiles à utiliser en pratique : il est difficile d'y trouver et
d'y rédiger des preuves.

\end{document}
